% GENERATED FILE. Edit canonical lessons, then run npm run generate:books.
% canonical-source-hash: fnv1a64:d7282e3c5df462fa
% canonical-lessons: JA-C105-hakama, JA-C105-yukata, JA-C105-honoo, JA-C105-akari, JA-C105-hikari

\chapter{Hakama, Cotton kimono, Flame, Lamp, Light}
\label{ch:ja-hakama-cotton-kimono-flame-lamp-light}
\hlchaptermodality{105}

\begin{chapteropening}
\textbf{By the end of this chapter:} \emph{I can say a hakama, a cotton kimono, a flame, a lamp and light.}

Complete the last lesson of chapter 105: I can say a hakama, a cotton kimono, a flame, a lamp and light.
\end{chapteropening}

\section[\texorpdfstring{\ja{はかま} (hakama)}{はかま (hakama)}]{\ja{はかま} (hakama) — a hakama (pleated trousers)}

\label{lesson:JA-C105-hakama}



\begin{quote}
Before the new one: say the Japanese for a ball, then the Japanese for a bell.
\end{quote}

\subsection*{You'll want to know: \ja{はかま}}
\textbf{\ja{はかま}} — \emph{hakama} — \textquotedblleft{}a hakama (pleated trousers)\textquotedblright{}.

\begin{grammarlens}[title={What you've built}]
One more word for this chapter.
\end{grammarlens}

\subsection*{Your turn}
\begin{itemize}
\raggedright
  \item \emph{Say it:} \emph{hakama}
  \item \emph{Say it:} \emph{hakama}, once more
  \item \emph{Say it:} say \emph{hakama}
  \item \emph{Recall:} say the Japanese for a ball, then the Japanese for a bell, then say \emph{hakama} again
\end{itemize}

\begin{tcolorbox}[breakable,colback=teal!4,colframe=teal!35!black,title={Before you move on}]
What does \ja{はかま} mean? (\textquotedblleft{}A hakama (pleated trousers)\textquotedblright{}.) Say it once more.
\end{tcolorbox}
\section[\texorpdfstring{\ja{ゆかた} (yukata)}{ゆかた (yukata)}]{\ja{ゆかた} (yukata) — a cotton kimono}

\label{lesson:JA-C105-yukata}



\begin{quote}
Before the new one: say the Japanese for a bell, then the Japanese for a hakama.
\end{quote}

\subsection*{You'll want to know: \ja{ゆかた}}
\textbf{\ja{ゆかた}} — \emph{yukata} — \textquotedblleft{}a cotton kimono\textquotedblright{}.

\begin{grammarlens}[title={What you've built}]
One more word for this chapter.
\end{grammarlens}

\subsection*{Your turn}
\begin{itemize}
\raggedright
  \item \emph{Say it:} \emph{yukata}
  \item \emph{Say it:} \emph{yukata}, once more
  \item \emph{Say it:} say \emph{yukata}
  \item \emph{Recall:} say the Japanese for a bell, then the Japanese for a hakama, then say \emph{yukata} again
\end{itemize}

\begin{tcolorbox}[breakable,colback=teal!4,colframe=teal!35!black,title={Before you move on}]
What does \ja{ゆかた} mean? (\textquotedblleft{}A cotton kimono\textquotedblright{}.) Say it once more.
\end{tcolorbox}
\section[\texorpdfstring{\ja{ほのお} (honoo)}{ほのお (honoo)}]{\ja{ほのお} (honoo) — a flame}

\label{lesson:JA-C105-honoo}



\begin{quote}
Before the new one: say the Japanese for a hakama, then the Japanese for a cotton kimono.
\end{quote}

\subsection*{You'll want to know: \ja{ほのお}}
\textbf{\ja{ほのお}} — \emph{honoo} — \textquotedblleft{}a flame\textquotedblright{}.

\begin{grammarlens}[title={What you've built}]
One more word for this chapter.
\end{grammarlens}

\subsection*{Your turn}
\begin{itemize}
\raggedright
  \item \emph{Say it:} \emph{honoo}
  \item \emph{Say it:} \emph{honoo}, once more
  \item \emph{Say it:} say \emph{honoo}
  \item \emph{Recall:} say the Japanese for a hakama, then the Japanese for a cotton kimono, then say \emph{honoo} again
\end{itemize}

\begin{tcolorbox}[breakable,colback=teal!4,colframe=teal!35!black,title={Before you move on}]
What does \ja{ほのお} mean? (\textquotedblleft{}A flame\textquotedblright{}.) Say it once more.
\end{tcolorbox}
\section[\texorpdfstring{\ja{あかり} (akari)}{あかり (akari)}]{\ja{あかり} (akari) — a lamp, a light}

\label{lesson:JA-C105-akari}



\begin{quote}
Before the new one: say the Japanese for a cotton kimono, then the Japanese for a flame.
\end{quote}

\subsection*{You'll want to know: \ja{あかり}}
\textbf{\ja{あかり}} — \emph{akari} — \textquotedblleft{}a lamp, a light\textquotedblright{}.

\begin{grammarlens}[title={What you've built}]
One more word for this chapter.
\end{grammarlens}

\subsection*{Your turn}
\begin{itemize}
\raggedright
  \item \emph{Say it:} \emph{akari}
  \item \emph{Say it:} \emph{akari}, once more
  \item \emph{Say it:} say \emph{akari}
  \item \emph{Recall:} say the Japanese for a cotton kimono, then the Japanese for a flame, then say \emph{akari} again
\end{itemize}

\begin{tcolorbox}[breakable,colback=teal!4,colframe=teal!35!black,title={Before you move on}]
What does \ja{あかり} mean? (\textquotedblleft{}A lamp, a light\textquotedblright{}.) Say it once more.
\end{tcolorbox}
\section[\texorpdfstring{\ja{ひかり} (hikari)}{ひかり (hikari)}]{\ja{ひかり} (hikari) — light}

\label{lesson:JA-C105-hikari}



\begin{quote}
Before the new one: say the Japanese for a flame, then the Japanese for a lamp.
\end{quote}

\subsection*{You'll want to know: \ja{ひかり}}
\textbf{\ja{ひかり}} — \emph{hikari} — \textquotedblleft{}light\textquotedblright{}.

\begin{grammarlens}[title={What you've built}]
One more word for this chapter.
\end{grammarlens}

\subsection*{Your turn}
\begin{itemize}
\raggedright
  \item \emph{Say it:} \emph{hikari}
  \item \emph{Say it:} \emph{hikari}, once more
  \item \emph{Say it:} say \emph{hikari}
  \item \emph{Recall:} say the Japanese for a flame, then the Japanese for a lamp, then say \emph{hikari} again
\end{itemize}

\begin{tcolorbox}[breakable,colback=teal!4,colframe=teal!35!black,title={Before you move on}]
What does \ja{ひかり} mean? (\textquotedblleft{}Light\textquotedblright{}.) Say it once more.
\end{tcolorbox}
